\documentclass[10pt, a4paper]{article}

\usepackage[
    ignoreheadfoot,
    top=0.8 in,
    bottom=0.8 in,
    left=0.75 in,
    right=0.75 in,
    footskip=0.4 in,
]{geometry}
\usepackage{titlesec}
\usepackage{tabularx}
\usepackage{array}
\usepackage[dvipsnames]{xcolor}
\definecolor{primaryColor}{RGB}{0, 0, 0}
\usepackage{enumitem}
\usepackage{fontawesome5}
\usepackage{amsmath}
\usepackage[
    pdftitle={Currículo — Leonardo Basso},
    pdfauthor={Leonardo Basso},
    colorlinks=true,
    urlcolor=primaryColor
]{hyperref}
\usepackage[pscoord]{eso-pic}
\usepackage{calc}
\usepackage{bookmark}
\usepackage{lastpage}
\usepackage{changepage}
\usepackage{paracol}
\usepackage{ifthen}
\usepackage{needspace}
\usepackage{iftex}
\usepackage{multicol}

\ifPDFTeX
    \input{glyphtounicode}
    \pdfgentounicode=1
    \usepackage[T1]{fontenc}
    \usepackage[utf8]{inputenc}
    \usepackage{lmodern}
\fi

\usepackage{charter}

\raggedright
\AtBeginEnvironment{adjustwidth}{\partopsep0pt}
\pagestyle{empty}
\setcounter{secnumdepth}{0}
\setlength{\parindent}{0pt}
\setlength{\topskip}{0pt}
\setlength{\columnsep}{0.15cm}
\pagenumbering{gobble}

\titleformat{\section}{\needspace{4\baselineskip}\bfseries\large}{}{0pt}{}[\vspace{1pt}\titlerule]

\titlespacing{\section}{
    -1pt
}{
    0.3 cm
}{
    0.2 cm
}

\renewcommand\labelitemi{$\vcenter{\hbox{\small$\bullet$}}$}
\newenvironment{highlights}{
    \begin{itemize}[
        topsep=0.10 cm,
        parsep=0.10 cm,
        partopsep=0pt,
        itemsep=0pt,
        leftmargin=0 cm + 10pt
    ]
}{
    \end{itemize}
}

\newenvironment{onecolentry}{
    \begin{adjustwidth}{
        0 cm + 0.00001 cm
    }{
        0 cm + 0.00001 cm
    }
}{
    \end{adjustwidth}
}

\newenvironment{twocolentry}[2][]{
    \onecolentry
    \def\secondColumn{#2}
    \setcolumnwidth{\fill, 4.5 cm}
    \begin{paracol}{2}
}{
    \switchcolumn \raggedleft \secondColumn
    \end{paracol}
    \endonecolentry
}

\begin{document}

    \begin{center}
        {\fontsize{25pt}{25pt}\selectfont \textbf{Leonardo Basso}}

        \vspace{5pt}

        \normalsize
        \href{mailto:leojunioyuri@gmail.com}{leojunioyuri@gmail.com}
        \quad|\quad
        \href{tel:+5535998794471}{+55 35 99879 4471}
        \quad|\quad
        \href{https://linkedin.com/in/leojunioyuri}{linkedin.com/in/leojunioyuri}
        \quad|\quad
        \href{https://github.com/leojunioyuri}{github.com/leojunioyuri}
        \quad|\quad
        \href{https://leojunioyuri.github.io}{leojunioyuri.github.io}
    \end{center}

    \vspace{5pt}

    % ─── DESCRIÇÃO ───────────────────────────────────────────────────────────
    \section{Descrição}
    \vspace{0.2 cm}
    \begin{onecolentry}
        Engenheiro de Software e Product Engineer com 4+ anos de experiência, formado em Ciência da Computação pela UFLA. Na Brendi, lidero a área de aquisição de consumidores, onde desenvolvi o MetaAds — sistema de tráfego pago autônomo com IA que levou de zero a 2.000+ restaurantes atendidos, com R\$20M+ em vendas rastreadas e retorno médio de 10$\times$ sobre o investimento.
    \end{onecolentry}

    % ─── EXPERIÊNCIA ─────────────────────────────────────────────────────────
    \section{Experiência}

    \begin{twocolentry}{Out 2024 -- Presente}
        \textbf{Product Engineer · Founding Team}, Brendi -- São José dos Campos, SP
    \end{twocolentry}
    \vspace{0.10 cm}
    \begin{onecolentry}
        \begin{highlights}
            \item Lidero o MetaAds: sistema end-to-end de tráfego pago automatizado com IA para 2.000+ restaurantes atendidos (600+ ativos). R\$20M+ em vendas rastreadas e retorno médio de 10$\times$ sobre o investimento.
            \item Arquitetura serverless em TypeScript/Node.js/Vue.js + GCP Cloud Functions, com migração Firestore $\rightarrow$ PostgreSQL/AlloyDB (220+ migrations Prisma, zero downtime).
            \item Anderson: agente de IA com ciclo semanal (Observar $\rightarrow$ Deliberar $\rightarrow$ Validar $\rightarrow$ Aplicar $\rightarrow$ Reportar), memória pgvector, loop de atribuição e 7 guardrails de segurança.
            \item Financial ledger próprio: R\$500k+/mês em verba de clientes gerida de forma autônoma via PIX + Meta API, sem operador humano no loop.
            \item PostHog (feature flags), Sentry (monitoramento, gauge metrics, 10\% trace sampling), Metabase com 150+ cards de analytics.
        \end{highlights}
    \end{onecolentry}

    \vspace{0.35 cm}

    \begin{twocolentry}{Ago 2023 -- Jul 2024}
        \textbf{Engenheiro de Software}, LEVTY -- Remoto
    \end{twocolentry}
    \vspace{0.10 cm}
    \begin{onecolentry}
        \begin{highlights}
            \item Especializado em plataformas all-in-one (BPM, ECM, CRM) para grandes organizações.
            \item Contribuiu para projeto de transição de grande organização educacional no Brasil, atingindo 100.000+ interações em 4 meses.
            \item Modelagem e implementação de processos lógicos para clientes enterprise; reuniões técnicas de alinhamento de requisitos.
        \end{highlights}
    \end{onecolentry}

    \vspace{0.35 cm}

    \begin{twocolentry}{Mar 2021 -- Jan 2023}
        \textbf{Estagiário Full Stack}, dti digital -- Belo Horizonte, MG
    \end{twocolentry}
    \vspace{0.10 cm}
    \begin{onecolentry}
        \begin{highlights}
            \item Desenvolvimento em JavaScript, TypeScript, React, C\#, .NET, Java e Spring Boot.
            \item Participou da reformulação de plataforma integrada de gestão para o maior cliente de mineração do Brasil.
            \item Testes com Jest, React Testing Library e xUnit. CI/CD e micro-frontends no Azure Cloud.
        \end{highlights}
    \end{onecolentry}

    \vspace{0.35 cm}

    \begin{twocolentry}{Jun 2019 -- Jan 2021}
        \textbf{Tech Lead \& Desenvolvedor Frontend}, Comp Júnior -- UFLA, Lavras, MG
    \end{twocolentry}
    \vspace{0.10 cm}
    \begin{onecolentry}
        \begin{highlights}
            \item Tech Lead: distribuição de conhecimento técnico, saúde das equipes e cultura ágil.
            \item Desenvolvimento frontend em empresa júnior com clientes reais e metodologia Scrum.
        \end{highlights}
    \end{onecolentry}

    % ─── FORMAÇÃO ─────────────────────────────────────────────────────────────
    \section{Formação}

    \begin{twocolentry}{Ago 2018 -- Dez 2023}
        \textbf{Universidade Federal de Lavras (UFLA)}, Bacharel em Ciência da Computação
    \end{twocolentry}
    \vspace{0.10 cm}
    \begin{onecolentry}
        \begin{highlights}
            \item Presidente do CIMVEST (2021): transformou o grupo no 1º núcleo de educação financeira da UFLA, crescendo de 10 para 40 membros. 1º Simpósio com 300+ espectadores.
        \end{highlights}
    \end{onecolentry}

    % ─── HABILIDADES ─────────────────────────────────────────────────────────
    \section{Habilidades}
    \vspace{0.1 cm}
    \begin{onecolentry}
        \begin{highlights}
            \item \textbf{Linguagens:} TypeScript, JavaScript, C\#, Node.js
            \item \textbf{Frontend:} Vue.js, Nuxt, React, Next.js, Tailwind CSS, Astro
            \item \textbf{Backend / Cloud:} GCP, Cloud Functions, Cloud Tasks, Firebase, Docker, GitHub Actions
            \item \textbf{Banco de Dados:} PostgreSQL, AlloyDB, MySQL, Supabase, Firebase, pgvector
            \item \textbf{IA \& Integrações:} OpenAI / LLMs, Meta Marketing API, WhatsApp API, Starkbank API
            \item \textbf{Observabilidade:} PostHog, Sentry, Metabase, Replit
            \item \textbf{Ferramentas:} Prisma, Zod, Biome, Git, Scrum / Metodologias Ágeis
            \item \textbf{Idiomas:} Português (nativo), Inglês (avançado / full professional), Italiano (elementar), Espanhol (básico)
        \end{highlights}
    \end{onecolentry}

    % ─── PROJETOS ─────────────────────────────────────────────────────────────
    \section{Projetos}

    \begin{twocolentry}{2023}
        \textbf{Pensão Mansão Azul} -- \href{https://www.pensaomansaoazul.com.br/}{pensaomansaoazul.com.br}
    \end{twocolentry}
    \vspace{0.10 cm}
    \begin{onecolentry}
        Site em React, TypeScript, Next.js e .NET para pensão estudantil em Lavras, MG. Atingiu $\sim$3.000 pessoas na primeira semana.
    \end{onecolentry}

    \vspace{0.25 cm}

    \begin{twocolentry}{2024}
        \textbf{Goal Invest} -- \href{https://goal-invest.vercel.app/}{goal-invest.vercel.app}
    \end{twocolentry}
    \vspace{0.10 cm}
    \begin{onecolentry}
        PWA para calcular rentabilidade de investimentos com aportes iniciais, contribuições mensais e metas financeiras.
    \end{onecolentry}

    \vspace{0.25 cm}

    \begin{twocolentry}{2017}
        \textbf{Beta Gun} -- \href{https://leojunioyuri.itch.io/beta-gun}{leojunioyuri.itch.io/beta-gun}
    \end{twocolentry}
    \vspace{0.10 cm}
    \begin{onecolentry}
        Jogo 2D de plataforma/shooter feito durante Game Jam na UFLA, em equipe de 5, usando Construct engine. Disponível gratuitamente no itch.io.
    \end{onecolentry}

\end{document}
